\documentclass[11pt,a4paper]{article}

\usepackage{amsmath}
\usepackage{amssymb}
\usepackage{booktabs}
\usepackage{multirow}
\usepackage{xcolor}
\usepackage{hyperref}
\usepackage{microtype}
\usepackage{geometry}
\usepackage{parskip}

\geometry{margin=2.5cm}

\definecolor{darkblue}{RGB}{0,50,120}
\hypersetup{colorlinks=true, linkcolor=darkblue, citecolor=darkblue, urlcolor=darkblue}

\title{\textbf{TrimAnalyser: Empirical Observations on Proof Structure\\
Across Subgraph Isomorphism Benchmark Families}}
\author{Arthur Gontier}
\date{June 2026 --- working notes}

\begin{document}

\maketitle

\begin{abstract}
We report empirical observations on the structure of VeriPB unsatisfiability certificates produced by
the Glasgow subgraph isomorphism solver across four benchmark families: LV, biochemical reactions
(bio), images-CVIU11, and meshes-CVIU11.
Each family yields a distinct \emph{proof fingerprint} --- a characteristic combination of step-type
mix, proof depth, algebraic complexity, and cone shape --- that traces directly back to the structural
properties of the underlying graphs.
We document the fingerprints quantitatively, explain their graph-theoretic origins, and derive
implications for family-aware \emph{solver} heuristics.
Grim is treated as the established trimming heuristic throughout; the open questions are at the
solver level.
\end{abstract}

\tableofcontents
\bigskip\hrule\bigskip

% ─────────────────────────────────────────────
\section{Setting and Methodology}

We analyse pseudo-Boolean proofs of unsatisfiability produced by the Glasgow subgraph isomorphism
solver~\cite{glasgow-icgt2020} extended with VeriPB proof
logging~\cite{sip-meets-cp-ijcai2020,e2e-verification-aaai2024}.
Proofs are trimmed by TrimAnalyser using the \textbf{Grim} heuristic, which has been established
as the best trimming strategy: during RUP verification, constraints already in the cone are served
first from a priority queue, steering conflict resolution toward the existing cone and minimising
cone growth.
Metrics are aggregated per instance by \texttt{scripts/aggregate\_results.jl} and joined with
per-graph structural features from \texttt{scripts/graph\_features.jl}.
All statistics below are averages over instances for which a valid trimmed cone was produced.

\paragraph{Cluster run.} Results are from a cluster run (June 2026) with flags:
\begin{quote}\ttfamily
./trimnalyser --threads 92,1 solve resolv verif allgraphs st=180 tt=6000 maxmem=50
\end{quote}
(solver timeout 180\,s, trimmer timeout 6\,000\,s, memory limit 50\,GB per thread,
no node-count filter).
Instance counts and proof coverage are in Table~\ref{tab:instance-counts}.

\begin{table}[h]
\centering
\begin{tabular}{lrrr}
\toprule
Family & Total & With cone & Coverage \\
\midrule
LV            & 4\,164 & 2\,486 & 60\% \\
bio           & 6\,711 & 5\,504 & 82\% \\
images-CVIU11 & 2\,997 &   506  & 17\% \\
meshes-CVIU11 & 2\,253 & 1\,730 & 77\% \\
\bottomrule
\end{tabular}
\caption{Instance counts. \emph{With cone}: instances for which TrimAnalyser produced a trimmed proof.
Instances without a cone fall into one of three categories: SAT (subgraph found --- no proof generated,
recorded in the \texttt{is\_sat} column of \texttt{cluster\_results.csv}), solver timeout (st=180\,s),
or memory-out (maxmem=50\,GB).
The 17\% coverage for images-CVIU11 likely reflects a high proportion of SAT instances and/or
timeouts; the exact split is available in the \texttt{is\_sat} and \texttt{skip\_reason} columns.}
\label{tab:instance-counts}
\end{table}

% ─────────────────────────────────────────────
\section{VeriPB Proof Step Types}
\label{sec:proof-steps}

Glasgow outputs proofs in the VeriPB pseudo-Boolean format~\cite{veripb-aaai2020}.
Each proof step derives a new linear inequality from existing ones.
TrimAnalyser distinguishes the following rule types in SIP proofs.

\begin{description}
  \item[POL (linear combination).] The general rule: derives a new inequality as a non-negative
    integer linear combination of existing constraints.
    Encoded in the proof as \texttt{pol} with a list of (coefficient, constraint-id) antecedents.
    The antecedent count per step (\texttt{pol\_ante\_mean}) measures algebraic complexity.

  \item[IA (implied and add).] A VeriPB step that asserts a new constraint by adding a referenced
    existing constraint to a supplied linear inequality.
    In Glasgow SIP proofs, IA appears exclusively as part of a three-step pattern:
    \begin{quote}\ttfamily
      pol $c_1\ c_2\ \ldots$ ;\ \ \emph{\% derives transient @p}\\
      ia @p $a x + b y \ge c$ ;\ \ \emph{\% adds @p to the explicit ineq, producing result}\\
      del @p\ \ \emph{\% removes the transient}
    \end{quote}
    The POL step derives an intermediate constraint \texttt{@p}; the IA step adds \texttt{@p} to
    the supplied inequality to produce the domain-reduction conclusion; \texttt{del} immediately
    discards the transient.
    Every IA step in the cone therefore has a paired POL step at the same depth, which explains
    why families with high IA fraction also show an approximately equal POL count
    (the paired-POL steps are distinct from the free-standing POL steps that appear in
    flat algebraic certificates).

  \item[RUP (reverse unit propagation).] The proof does not supply the derivation explicitly;
    the verifier reconstructs it by unit propagation.
    RUP steps are cheap to write but more expensive to verify.

  \item[RED (redundance-based extension).] Introduces an auxiliary constraint that is redundant
    with respect to the proof goal.
    RED does \emph{not} appear in SIP proofs: SIP is a pure feasibility problem with no
    optimisation objective, so redundance-based reasoning is never required.
    It is listed here for completeness.
\end{description}

% ─────────────────────────────────────────────
\section{Benchmark Families}

\subsection{LV (Larrosa--Valiente)}
\label{sec:lv}

The LV benchmark was introduced by Larrosa and Valiente~\cite{larrosa-valiente-2002} as a
collection of 112 graphs spanning a range of structural types: connected, biconnected,
triconnected, bipartite, planar, and regular, with node counts ranging from 10 to 6\,671.
Instances are formed by taking all pairs $(P, T)$ with $|V(P)| \le |V(T)|$ from this pool.
The benchmark is hosted as part of the newSIP benchmark suite of Solnon~\cite{newsip-gbrpr2019}.

\paragraph{Graph structure.} LV graphs are among the densest in the benchmark suite.
Pattern graphs average 259 nodes and 2\,836 edges (density 0.168, mean degree 13.7).
The degree sequence is highly heterogeneous (degree variance 92.4): the pool contains both
sparse and dense graph types, and density varies widely across pairs.
About 27\% of LV pattern graphs are bipartite and 15\% are regular.
The average pattern-to-target node ratio is 0.30 and the density ratio is 42.5,
meaning patterns are typically much denser than their targets.
The degree compatibility fraction (\texttt{degree\_compat\_frac}) is 0.810: in roughly one
in five LV pairs, some pattern vertex has degree exceeding the target's maximum degree, making
the instance trivially UNSAT from degree alone.

\subsection{Bio (Biochemical Reactions)}
\label{sec:bio}

The bio benchmark consists of directed graphs representing biochemical reaction networks
extracted from the BioModels repository~\cite{bonnici-bioinformatics2013}.
Graphs encode metabolic networks: one vertex class represents reactions, the other represents
chemical species, and edges connect species to the reactions they participate in.
Instances are formed by taking all $(P, T)$ pairs from a pool of 136 graphs with 9--386 nodes.
The benchmark is included in the newSIP suite~\cite{newsip-gbrpr2019}.

\paragraph{Graph structure.} Bio patterns are extremely sparse (density 0.086, mean degree 1.61)
with no triangles (clustering coefficient 0.000), reflecting the bipartite structure of
reaction networks.
Diameters are long: 9.4 for patterns and 12.6 for targets on average.
The low degree and absence of triangles means the graphs are locally tree-like, with the embedding
constraint transmitted along long paths.
About 56\% of pattern graphs pass an undirected bipartiteness test.

\paragraph{\textbf{TODO:} bipartite structure not exploited.}
Biochemical reaction networks are bipartite \emph{by construction} (reaction nodes vs.\
metabolite nodes), but the LAD files are treated as general undirected graphs throughout this
analysis.
It is an open question whether exploiting bipartite structure in the solver
(e.g.\ restricting search to bipartite-structure-preserving mappings) or in the trimmer
would significantly change the proof fingerprint or improve performance.
This should be investigated as a dedicated experiment.

\subsection{Images-CVIU11}
\label{sec:images}

The images-CVIU11 family originates from Damiand~et~al.~\cite{damiand-solnon-cviu2011}, a study of
submap isomorphism algorithms for image analysis.
Pattern graphs are derived from small segmented image regions (22--151 nodes); target graphs are
region adjacency graphs (RAGs) of full images (1\,072--5\,972 nodes).
Each node represents an image region and each edge represents spatial adjacency.
Instances are formed by crossing all patterns against all targets with $|V(P)| \le |V(T)|$
(44 patterns $\times$ 146 targets).

\paragraph{Graph structure.} Image RAGs are sparse and irregular.
Pattern graphs average 83 nodes with density 0.053 and mean degree 2.78; triangles are rare
(clustering coefficient 0.11).
Diameters are long (17.4 for patterns; target diameters are often undefined due to disconnected
components from image segmentation).
The most striking structural feature is the size disparity: the average pattern-to-target node
ratio is only 0.026 (pattern $\approx$ 82 nodes, target $\approx$ 3\,696 nodes).
Despite both being sparse in absolute terms, patterns are 67$\times$ denser than their targets.
Degree compatibility is trivially satisfied (fraction 1.000) since target graphs are large enough
to accommodate any pattern degree.

\subsection{Meshes-CVIU11}
\label{sec:meshes}

Also from Damiand~et~al.~\cite{damiand-solnon-cviu2011}, meshes-CVIU11 contains graphs derived
from 3D mesh combinatorial maps of three-dimensional objects.
Nodes correspond to mesh facets or vertices; edges reflect topological adjacency in the mesh.
The family has 6 patterns (40--199 nodes) and 503 targets (208--5\,873 nodes),
with instances formed by crossing all pairs.

\paragraph{Graph structure.} Mesh graphs exhibit the regularity characteristic of triangulated
3D surfaces: bounded degree (mean degree 5.33 for patterns), high local clustering (0.325),
and relatively uniform degree sequences (degree variance 19.3 for patterns).
Unlike images, mesh targets have high degree variance (1\,312) because large meshes can include
both high-valence interior vertices and low-degree boundary vertices.
Pattern-to-target node ratio is 0.117 and degree compatibility fraction is 0.990.

% ─────────────────────────────────────────────
\section{Proof Fingerprints}

Table~\ref{tab:proof-features} summarises the key proof-structural metrics per family.
We discuss each fingerprint below.

\begin{table}[h]
\centering
\renewcommand{\arraystretch}{1.15}
\begin{tabular}{lrrrr}
\toprule
Metric & LV & bio & images & meshes \\
\midrule
\multicolumn{5}{l}{\emph{Proof size (cone steps)}} \\
\quad Total cone (OPB + PBP)    & 50\,480 & 8\,473 & 229\,147 & 113\,996 \\
\quad OPB cone (original)       & 43\,415 &  2\,220 &  82\,249 & 101\,049 \\
\quad PBP cone (derived)        &  7\,065 &  6\,253 & 146\,898 &  12\,947 \\
\midrule
\multicolumn{5}{l}{\emph{Step-type fractions}} \\
\quad $f_{\mathrm{POL}}$   & 0.934 & 0.719 & 0.496 & 0.998 \\
\quad $f_{\mathrm{IA}}$    & 0.037 & 0.268 & 0.503 & 0.002 \\
\quad $f_{\mathrm{RUP}}$   & 0.029 & 0.013 & 0.001 & 0.001 \\
\midrule
\multicolumn{5}{l}{\emph{Proof depth}} \\
\quad Max depth              &   9.2 &   6.2 &  93.6 &   2.0 \\
\quad Mean depth             &   3.3 &   1.6 &   2.9 &   1.0 \\
\quad Depth entropy          &  0.29 &  1.09 &  2.18 &  0.01 \\
\quad Bottom fraction        &  0.97 &  0.89 &  0.66 &  1.00 \\
\midrule
\multicolumn{5}{l}{\emph{Algebraic complexity (POL steps)}} \\
\quad Antecedents per POL    & 34.96 &  9.06 &  5.45 & 15.97 \\
\quad OPB fraction of antes  &  0.977 & 0.931 & 0.830 & 1.000 \\
\quad POL unlocks RUP burst  & 5.5\% & 49.7\% & 72.3\% & 0.3\% \\
\midrule
\multicolumn{5}{l}{\emph{Resolv loop}} \\
\quad Pattern shrinkage      &  0.495 &  0.353 &  0.129 &  0.377 \\
\quad Iterations             &  0.294 &  0.387 &  0.074 &  0.406 \\
\midrule
\multicolumn{5}{l}{\emph{Cone shape}} \\
\quad Max width              & 4\,618 & 2\,836 & 47\,776 & 12\,903 \\
\quad Width CV               &  1.46  &  1.43  &   2.57  &   1.41  \\
\quad Literal weakening rate &  0.32  &  0.37  &  0.28   &  0.45   \\
\midrule
\multicolumn{5}{l}{\emph{Runtime}} \\
\quad Avg runtime (ms)   & 20\,777 & 518 & 65\,533 & 22\,046 \\
\bottomrule
\end{tabular}
\caption{Proof-structural metrics averaged over instances with valid cone data.
\emph{Bottom fraction}: fraction of cone steps at depth $\le 2$.
\emph{POL unlocks RUP burst}: fraction of instances where a depth band contains a POL step
followed at depth$+1$ by a RUP count $>5\times$ the per-depth mean.
\emph{Literal weakening rate}: $(L_{\mathrm{cone}} - L_{\mathrm{smol}}) / L_{\mathrm{cone}}$.}
\label{tab:proof-features}
\end{table}

\subsection{LV: Large-Scale Algebraic Certificates from Dense Graphs}
\label{sec:fingerprint-lv}

LV proofs are \textbf{algebraically heavy}: the average POL step combines 34.96 antecedents
(the highest of all families), 97.7\% of which are original OPB constraints.
The step-type mix is nearly pure POL ($f_{\mathrm{POL}} = 0.934$) with negligible IA
and RUP fractions, indicating almost no domain-reduction chains during search.
Proofs are \textbf{shallow} (max depth 9.2, entropy 0.29, bottom fraction 0.97): almost all steps
are within two hops of the axioms.

The resolv loop is \textbf{highly effective}: pattern graphs shrink by 49.5\% on average.
The PBP cone is small relative to the OPB cone (7\,065 vs.\ 43\,415): the certificate is
essentially a short sequence of large algebraic aggregations of many original constraints.

\paragraph{Interpretation.} Dense, heterogeneous graphs create many simultaneously tight degree
and adjacency constraints.
When a pattern vertex has higher degree than any target vertex
(degree\_compat\_frac $= 0.810$), the infeasibility is immediate from degree algebra.
In less trivial cases, the solver combines many neighbourhood constraints in one large POL step.
The resulting certificate is wide (many antecedents) but shallow.

\subsection{Bio: Mixed Propagation--Algebra on Sparse Networks}
\label{sec:fingerprint-bio}

Bio proofs show the most complex \textbf{step-type interleaving}: $f_{\mathrm{POL}} = 0.719$ and
$f_{\mathrm{IA}} = 0.268$.
Depth entropy is the highest of all families (1.094), meaning proof steps are spread across
many depth levels with no dominant tier.
Almost half of bio instances (49.7\%) exhibit the ``POL unlocks RUP burst'' pattern.

The resolv loop is moderately effective (0.353 pattern shrinkage, mean 0.387 non-zero iterations),
suggesting the UNSAT core emerges gradually.

\paragraph{Interpretation.} Tree-like, long-diameter networks require the solver to propagate
domain reductions along long paths before the infeasibility becomes apparent.
The IA steps (paired with transient POL steps, as in Section~\ref{sec:proof-steps}) chain the
resulting domain reductions along the path until a contradiction accumulates.
The high depth entropy reflects the long diameter: there is no preferred depth at which most
reasoning happens.

\subsection{Images-CVIU11: Propagation Cascades on Irregular Adjacency Graphs}
\label{sec:fingerprint-images}

Images-CVIU11 proofs are the most extreme in the dataset along multiple axes.
Step types split almost equally between POL and IA ($f_{\mathrm{POL}} \approx f_{\mathrm{IA}}
\approx 0.50$); since every IA step is paired with a transient POL step
(Section~\ref{sec:proof-steps}), approximately \emph{half of all proof steps} are
domain-reduction triplets \texttt{pol\,;\,ia\,;\,del}.
Proofs are \textbf{very deep}: the maximum depth averages 93.6 and entropy reaches 2.18.
The \textbf{PBP cone dominates}: 146\,898 derived steps vs.\ 82\,249 original constraints
(ratio 1.79, reversed compared to all other families).
The ``POL unlocks RUP burst'' pattern occurs in 72.3\% of instances.
The resolv loop is the \textbf{least effective} (pattern shrinkage 0.129).

\paragraph{Interpretation.} The 0.026 pattern-to-target node ratio means the solver eliminates
candidate mappings over a 37$\times$ larger target, generating many \texttt{pol\,;\,ia\,;\,del}
triplets.
These chain into a deep proof DAG: algebraic steps derive new bounds, which unlock cascades of
domain reductions, which force the next algebraic step.
Because the conflict spans the full irregular adjacency structure of the pattern, the resolv loop
has little room to shrink it.

\subsection{Meshes-CVIU11: Flat One-Wave Certificates from Regular 3D Topology}
\label{sec:fingerprint-meshes}

Meshes-CVIU11 is the most algebraically uniform family.
Proofs are \textbf{almost perfectly flat}: mean depth 1.002, depth entropy 0.011, bottom
fraction 1.000.
The step-type mix is pure POL ($f_{\mathrm{POL}} = 0.998$, $f_{\mathrm{IA}} = 0.002$),
meaning almost no domain-reduction triplets occur.
Every POL antecedent is an original OPB constraint (\texttt{pol\_opb\_frac} $= 1.000$).
The resolv loop achieves moderate pattern shrinkage (0.377) in few iterations.

\paragraph{Interpretation.} High local clustering (0.325) means triangles enforce three mutual
constraints simultaneously, making the combined constraint tight enough to certify UNSAT
algebraically without any propagation chain.
The resolv loop is effective because the topological obstruction is localised to a compact
substructure of the pattern.

% ─────────────────────────────────────────────
\section{Structural Drivers: Graph Properties $\to$ Proof Shape}
\label{sec:drivers}

Table~\ref{tab:graph-features} collects the graph-structural features that best predict the
observed proof fingerprints.

\begin{table}[h]
\centering
\renewcommand{\arraystretch}{1.15}
\begin{tabular}{lrrrr}
\toprule
Metric & LV & bio & images & meshes \\
\midrule
\multicolumn{5}{l}{\emph{Pattern graph}} \\
\quad Nodes                 & 259 &  29 &  83 & 111 \\
\quad Density               & 0.168 & 0.086 & 0.053 & 0.072 \\
\quad Mean degree           & 13.7 & 1.61 & 2.78 & 5.33 \\
\quad Degree variance       & 92.4 & 1.02 & 0.20 & 19.3 \\
\quad Clustering coeff.     & 0.275 & 0.000 & 0.109 & 0.325 \\
\quad Diameter              &  5.97 & 9.39 & 17.39 & 10.15 \\
\quad Bipartite (\%)        & 27\% & 56\% & 0\% & 0\% \\
\midrule
\multicolumn{5}{l}{\emph{Target graph}} \\
\quad Nodes                 & 1\,403 &  81 & 3\,696 & 1\,887 \\
\quad Density               & 0.057 & 0.038 & 0.001 & 0.005 \\
\quad Degree variance       & 248 & 1.77 & 0.13 & 1\,313 \\
\midrule
\multicolumn{5}{l}{\emph{Relational}} \\
\quad Node ratio            & 0.299 & 0.481 & 0.026 & 0.117 \\
\quad Density ratio         &  42.5 &  3.19 & 66.8 &  32.5 \\
\quad Degree compat.\ frac. & 0.810 & 0.985 & 1.000 & 0.991 \\
\bottomrule
\end{tabular}
\caption{Graph-structural features averaged over instances in the graph features dataset.
\emph{Density ratio}: pattern density / target density.
\emph{Degree compat.\ frac.}: fraction of pattern vertices with degree $\le \max_{v \in T}\deg(v)$.}
\label{tab:graph-features}
\end{table}

\paragraph{Width--depth tradeoff.}
A consistent inverse relationship holds between mean POL antecedents and proof depth:
\[
  \underbrace{\text{LV: 35 antes, depth 9.2}}_{\text{wide \& shallow}}
  \quad
  \underbrace{\text{meshes: 16 antes, depth 2.0}}_{\text{very flat}}
  \quad
  \underbrace{\text{bio: 9 antes, depth 6.2}}_{\text{balanced}}
  \quad
  \underbrace{\text{images: 5 antes, depth 93.6}}_{\text{narrow \& deep}}
\]
Families that aggregate many constraints per step produce wide but shallow proofs;
families that proceed via many small domain-reduction triplets produce narrow but deep proofs.

\paragraph{Graph density and algebraic step size.}
Higher pattern density correlates with larger \texttt{pol\_ante\_mean}.
Dense graphs (LV: 0.168) generate many simultaneously active neighbourhood constraints;
combining them requires a single large POL step.
Sparse graphs (images: 0.053) produce fewer simultaneous constraints per neighbourhood, so
each POL step is small and many domain-reduction triplets are needed instead.

\paragraph{Clustering and proof flatness.}
High clustering (meshes: 0.325) enables flat, OPB-heavy proofs.
Triangles impose three simultaneous constraints, making algebraic UNSAT certification
possible without any propagation chain.
Zero clustering (bio: 0.000) forces the solver into long propagation chains.

\paragraph{Diameter and IA fraction.}
Long diameter correlates with high IA fraction.
Bio (diameter 9.4) and images (diameter 17.4) have the highest IA fractions (26.8\% and 50.3\%).
Long propagation paths produce long chains of \texttt{pol\,;\,ia\,;\,del} triplets.

\paragraph{Node ratio and resolv effectiveness.}
The pattern-to-target node ratio is the strongest predictor of resolv loop effectiveness.
Images (ratio 0.026) achieve shrinkage 0.129; bio (ratio 0.481) achieves 0.353;
LV (ratio 0.299) achieves 0.495.
A small ratio means the conflict is spread across many candidate mappings, making it
hard for the resolv loop to isolate a compact UNSAT core.

% ─────────────────────────────────────────────
\section{Implications for Solver Heuristics}
\label{sec:solver-heuristics}

The Grim heuristic has been established as the best trimming strategy and is used uniformly.
The open research questions are at the \emph{solver} level.
The two primary axes are \textbf{variable ordering} (which pattern vertex to branch on next)
and \textbf{preprocessing / propagation ordering} (how aggressively to filter domains before
and during search).
A full A/B evaluation across these axes is planned for M4.

\paragraph{Mesh-like families: preprocessing first.}
The flat, OPB-heavy proof structure indicates that infeasibility is fully detectable from the
original degree and adjacency constraints before any deep search.
Aggressive preprocessing (degree-based domain filtering, neighbourhood compatibility checks,
triangle-aware propagation given clustering 0.325) should reduce the search space to near-zero.
Variable ordering plays a secondary role.

\paragraph{LV-like families: exploit trivial degree incompatibility early.}
With degree\_compat\_frac $= 0.810$, one in five LV instances can be certified UNSAT by a
lightweight pre-check (does any pattern vertex exceed the target's max degree?).
For non-trivial cases, the high \texttt{pol\_ante\_mean} (34.96) suggests that many constraints
are simultaneously tight at the conflict; a global constraint propagator covering all
neighbourhood constraints at once may be more effective than local pairwise propagation.

\paragraph{Image-like families: variable ordering is critical.}
The 0.026 node ratio means the initial candidate set for each pattern vertex is enormous.
Branching first on the most-constrained pattern vertex can prune large subtrees early.
The 72.3\% POL-before-RUP-burst signature means heuristics that trigger domain-reduction cascades
early in the search would reduce both solve time and proof depth.

\paragraph{Bio-like families: propagation along long paths.}
With mean degree 1.61 and diameter 9.4, constraint propagation must reach the far end of long
chains before detecting a contradiction.
Variable ordering based on graph topology (e.g.\ BFS-order from a high-betweenness seed vertex)
may guide propagation toward the conflict end of the path more efficiently.
Multi-iteration resolv should be prioritised here given moderate but consistent shrinkage (0.353)
across several iterations.

% ─────────────────────────────────────────────
\section{Open Questions}
\label{sec:open}

\begin{enumerate}
  \item \textbf{Remaining families.}
    The phase, scalefree, and si families have too few instances with cone data in the current run
    to characterise reliably.
    A targeted run with relaxed timeout or node bounds is needed.

  \item \textbf{Depth entropy as a classifier.}
    Depth entropy cleanly separates the four families (0.01, 0.29, 1.09, 2.18) and may serve
    as the primary axis of a two-dimensional taxonomy (entropy $\times$ ia\_frac) for runtime
    solver configuration selection.

  \item \textbf{Intra-family scaling laws.}
    Does proof size scale as $O(|V(P)|)$, $O(|V(T)|)$, or $O(|V(P)| \cdot |V(T)|)$ within each
    family?

  \item \textbf{Bio bipartite structure.}
    (See TODO in Section~\ref{sec:bio}.)
    A dedicated experiment exploiting bipartiteness in the solver and trimmer for bio instances.

  \item \textbf{Degree variance as a predictor within LV.}
    LV's internal diversity (bipartite, regular, planar sub-families) suggests that
    \texttt{pat\_deg\_var} and \texttt{pat\_is\_bipartite} stratify LV into distinct
    sub-fingerprints.

  \item \textbf{Images-CVIU11 coverage.}
    Determine the exact SAT / timeout / memout split for images to understand whether
    low coverage is structural or a resource limit.

  \item \textbf{Width--depth tradeoff: formal characterisation.}
    The observed inverse relationship between \texttt{pol\_ante\_mean} and
    \texttt{cone\_depth\_max} invites a proof-complexity explanation in terms of the Glasgow
    proof strategy.
\end{enumerate}

% ─────────────────────────────────────────────
\bibliographystyle{plain}
\bibliography{notes}

\end{document}
